\section*{Chapter \ref{ch:b2:eruptive-styles} --- Volcanic Eruptions}

\begin{solution}{exo:b2:eruptive-styles:1}
(a) $\qty{1.3}{km^3} = \qty{1.3e9}{m^3}$, between $10^9$ and $10^{10}$: VEI~5.
(b) $\qty{2.9e8}{m^3}$: VEI~4. (c) $\qty{2e6}{m^3}$, between $10^6$ and $10^7$:
VEI~2. (d) VEI~6 starts at \qty{e10}{m^3}, so the bulk \omterm{def:b2:eruptive-styles:tephra}{tephra} must be at least
\qty{10}{km^3}; ``more than \qty{5}{km^3}'' is a lower bound, not a contradiction.
\end{solution}

\begin{solution}{exo:b2:eruptive-styles:2}
(a) Ash (finer than \qty{2}{mm}). (b) A lapillus (\qtyrange{2}{64}{mm}).
(c) A block: larger than \qty{64}{mm}, angular, solid when ejected. (d) A
\omterm{def:b2:eruptive-styles:tephra}{volcanic bomb}, a spindle bomb: larger than \qty{64}{mm} and shaped while molten
in flight.
\end{solution}

\begin{solution}{exo:b2:eruptive-styles:3}
(a) $0.304\times(10^6)^{0.241} = 0.304\times27.9 = \qty{8.5}{km}$;
$0.304\times(10^8)^{0.241} = 0.304\times84.7 = \qty{25.8}{km}$.
(b) $Q = (35/0.304)^{1/0.241} = 115.1^{4.149} = \qty{3.6e8}{kg/s}$, and $I =
8.55 + 3 = 11.6$. (c) $Q \propto H^{4.15}$, so a height \qty{20}{\percent} too high
overestimates $Q$ by $1.2^{4.15} = 2.1$.
\end{solution}

\begin{solution}{exo:b2:eruptive-styles:4}
$V_{\mathrm{DRE}} = 0.29\times1000/2500 = \qty{0.116}{km^3}$; $m =
\num{0.29e9}\times1000 = \qty{2.9e11}{kg}$; $M = \log_{10}(\num{2.9e11}) - 7 = 11.46
- 7 = 4.46$.
\end{solution}

\begin{solution}{exo:b2:eruptive-styles:5}
$k = \ln(40/10)/(15 - 5) = 1.386/10 = \qty{0.139}{km^{-1}}$; $T_0 = 40\,e^{0.139\times5}
= 40\times2.00 = \qty{80}{cm}$; $\ln2/k = \qty{5.0}{km}$. With $T_0 =
\qty{8.0e-4}{km}$: $V = 2\times\num{8.0e-4}/0.139^2 = \qty{0.083}{km^3} =
\qty{8.3e7}{m^3}$, VEI~3.
\end{solution}

\begin{solution}{exo:b2:eruptive-styles:6}
(a) $\int_0^x T_0e^{-ky}2y\,dy = 2T_0\big[-\tfrac{y}{k}e^{-ky} - \tfrac{1}{k^2}e^{-ky}\big]_0^x
= \tfrac{2T_0}{k^2}[1 - (1 + kx)e^{-kx}]$. (b) Iterating $u = \ln(2(1 + u))$ from
$u = 1$: $1.386$, $1.564$, $1.635$, $1.662$, $1.672$, \dots, converging to
$u = 1.678$, for which $(1 + u)e^{-u} = 2.678\times0.1867 = 0.500$. (c) $x_{1/2} =
1.678/0.139 = \qty{12.1}{km}$, inside the outermost isopach ($\sqrt A =
\qty{15}{km}$); yet the volume beyond that isopach is $(1 + 2.09)e^{-2.09} =
\qty{38}{\percent}$ of the total. Almost two fifths of the deposit lies where
nothing was measured: the volume rests on the extrapolation.
\end{solution}

\begin{solution}{exo:b2:eruptive-styles:7}
(a) $p_s = (5.0/0.444)^2 = \qty{127}{MPa}$, at $127/0.025 = \qty{5.1}{km}$.
(b) $x = (5.0 - 0.444\sqrt{10})/100 = (5.0 - 1.404)/100 = 0.0360$;
$V_g = 0.0360\times8.314\times1123/(0.018015\times10^7) = \qty{1.86e-3}{m^3}$;
$V_m = 0.964/2300 = \qty{4.19e-4}{m^3}$; $\phi = 1.86/(1.86 + 0.419) = 0.82$.
(c) A basaltic melt is thousands to billions of times less viscous than a
rhyolitic one (\cref{ch:b2:magma-emplacement}): its bubbles rise,
merge and escape, and the gas leaves through a permeable froth or at the
vent, so the \omterm{def:b2:magma-emplacement:magma}{magma} degasses in an open system instead of keeping its gas; its
fragmentation, when it happens, is confined to a fountain near the vent. Slower
ascent leaves more time for this.
\end{solution}

\begin{solution}{exo:b2:eruptive-styles:8}
(a) $u_z(0) = 0.75\times10^7/(\pi\times(5000)^2) = \qty{0.095}{m} = \qty{95}{mm}$.
(b) $z_0 = 3.8/0.766 = \qty{5.0}{km}$. (c) $(1 + 1)^{-3/2} = 0.35$ of the
maximum.
\end{solution}

\begin{solution}{exo:b2:eruptive-styles:9}
The greatest range is at \ang{45}: $R = v^2/g = 100^2/9.81 = \qty{1.0}{km}$, and
$200^2/9.81 = \qty{4.1}{km}$. Air resistance shortens the flight, so these are
upper bounds, close to the truth for large blocks and far too long for small ones.
Since $R = v^2/g$, the hazard zone grows as the square of the launch speed:
doubling the speed of the explosion quadruples its radius.
\end{solution}

\begin{solution}{exo:b2:eruptive-styles:10}
(a) With $T = T_0e^{-k_1x}$ for $x < X$ and $T = T_Xe^{-k_2(x - X)}$ beyond, $T_X =
T_0e^{-k_1X}$ (continuity), use $\int_0^X ye^{-ky}dy = \frac{1}{k^2} - e^{-kX}\big(\frac{X}{k} +
\frac{1}{k^2}\big)$ and $\int_X^\infty ye^{-ky}dy = e^{-kX}\big(\frac{X}{k} +
\frac{1}{k^2}\big)$:
the proximal part gives $2T_0\big[\frac{1}{k_1^2} - e^{-k_1X}\frac{k_1X + 1}{k_1^2}\big]$ and the
distal part $2T_Xe^{k_2X}\int_X^\infty ye^{-k_2y}dy = 2T_X\frac{k_2X + 1}{k_2^2}$; with $T_X =
T_0e^{-k_1X}$ their sum is the formula of the method box.
(b) $T_0 = \qty{0.01285}{km}$: $2T_0/k_1^2 = 0.0257/0.1332 = \qty{0.193}{km^3}$;
$e^{-0.365\times8.56} = 0.0440$; $(0.138\times8.56 + 1)/0.138^2 = 114.5$ and
$(0.365\times8.56 + 1)/0.365^2 = 30.96$; the correction is $0.0257\times0.0440\times83.6
= \qty{0.094}{km^3}$ and $V = \qty{0.287}{km^3}$, the published \qty{0.29}{km^3}.
Beyond $X$ lies $2T_X(X/k_2 + 1/k_2^2) = 2\times\num{5.65e-4}\times114.5 =
\qty{0.129}{km^3}$, that is \qty{45}{\percent}.
\end{solution}

\begin{solution}{exo:b2:eruptive-styles:11}
(a) $100^{1/4} = 3.16$; the empirical laws give $100^{0.241} = 3.03$ and
$100^{0.226} = 2.83$. (b) $H \propto N^{-3/4}$: a more strongly stratified
atmosphere stops the column lower; for $N$ larger by \qty{20}{\percent},
$1.2^{-0.75} = 0.87$, a height lower by \qty{13}{\percent}. (c) Weak plumes are
bent over by the wind, which lowers them, while moist air releases latent heat in
them, which raises them; these effects weigh differently on small and large
eruptions and are not in the simple law. And the rates are known only to within
factors of several, an error in the variable on the horizontal axis that biases
a fitted slope towards zero.
\end{solution}

\begin{solution}{exo:b2:eruptive-styles:12}
(a) $132/400\times100 = 33$ per century before 1900; $95/125\times100 = 76$ per
century since. (b) $33/76 = \qty{43}{\percent}$. (c) $21/56 = \qty{38}{\percent}$:
the same under-recording, by a factor of two to three. (d) Larger eruptions are
seen from farther away and leave thicker, wider deposits that survive erosion
longer; a VEI~4 eruption on a remote island in 1700 could pass unrecorded, a
VEI~6 eruption hardly.
\end{solution}

\begin{solution}{pb:b2:eruptive-styles:1}
\textbf{1.} $\ln T = 5.298$, $4.605$, $3.871$, $2.639$, $1.386$. The slopes between
successive isopachs are $-0.0693$, $-0.0367$, $-0.0308$ and $-0.0313$ per km: one
steep segment between 10 and \qty{20}{km}, a gentle one from 40 to
\qty{120}{km}; the break lies between 20 and \qty{40}{km}.

\textbf{2.} $k_1 = \ln(200/100)/10 = \qty{0.0693}{km^{-1}}$; $T_0 = 200\,e^{0.693} =
\qty{400}{cm}$.

\textbf{3.} $167e^{-1.244} = \qty{48.1}{cm}$, $167e^{-2.488} = \qty{13.9}{cm}$,
$167e^{-3.732} = \qty{4.00}{cm}$: within \qty{1}{\percent} of 48, 14 and 4.0.

\textbf{4.} $400e^{-0.0693X} = 167e^{-0.0311X}$ gives $X = \ln(400/167)/(0.0693 -
0.0311) = 0.873/0.0382 = \qty{22.9}{km}$.

\textbf{5.} $\ln2/0.0693 = \qty{10.0}{km}$ and $\ln2/0.0311 = \qty{22.3}{km}$.

\textbf{6.} $T_0 = \qty{0.0040}{km}$: $V = 2\times0.0040/0.0693^2 =
\qty{1.67}{km^3}$.

\textbf{7.} $e^{-0.0693\times22.9} = 0.205$; $(0.0311\times22.9 + 1)/0.0311^2 = 1770$;
$(0.0693\times22.9 + 1)/0.0693^2 = 539$; correction $2\times0.0040\times0.205\times1231 =
\qty{2.02}{km^3}$; $V = 1.67 + 2.02 = \qty{3.68}{km^3}$, of which the correction is
\qty{55}{\percent}.

\textbf{8.} $3.68/1.67 = 2.2$. The fine distal ash is thin, easily eroded,
reworked or fallen into the sea; in an old deposit it is often lost or never
mapped, and the volume is then underestimated by a factor of two or more.

\textbf{9.} $\qty{3.68e9}{m^3}$ lies between $10^9$ and $10^{10}$: VEI~5.

\textbf{10.} $V_{\mathrm{DRE}} = 3.68\times900/2400 = \qty{1.38}{km^3}$.

\textbf{11.} $m = \num{3.68e9}\times900 = \qty{3.31e12}{kg}$.

\textbf{12.} $M = \log_{10}(\num{3.31e12}) - 7 = 12.52 - 7 = 5.52$.

\textbf{13.} $3.68/1.3 = 2.8$ times.

\textbf{14.} A column higher than \qty{25}{km}, and a Plinian style.

\textbf{15.} $Q = \num{3.31e12}/36000 = \qty{9.2e7}{kg/s}$; $I = 7.96 + 3 = 11.0$.

\textbf{16.} $0.304\times(\num{9.2e7})^{0.241} = 0.304\times83.0 = \qty{25.2}{km}$.

\textbf{17.} $0.345\times(\num{9.2e7})^{0.226} = 0.345\times63.1 = \qty{21.8}{km}$.

\textbf{18.} The measured \qty{31}{km} is \qty{23}{\percent} above the first
prediction and \qty{42}{\percent} above the second, within the
\qtyrange{53}{60}{\percent} scatter of the laws: consistent.

\textbf{19.} $Q = (31/0.304)^{1/0.241} = 102.0^{4.149} = \qty{2.2e8}{kg/s}$, 2.3 times
the mean rate. The clasts record the highest column, at the peak of the
eruption, while question~15 averages over the whole phase; and part of the mass
may be missing from the deposit (distal ash, \omterm{def:b2:eruptive-styles:pdc}{pyroclastic flows}).

\textbf{20.} Yes: 22 to \qty{31}{km}, well above the tropopause at \qty{16}{km}.
Gas carried into the stratosphere spreads around the globe; the key gas is
sulfur dioxide, and the twenty million tonnes that Pinatubo put there in 1991
cooled the globe by about half a degree for two years.

\textbf{21.} $\lambda_3 = 399/75 = 5.32$, $\lambda_4 = 95/125 = 0.760$, $\lambda_5 = 22/225 =
0.0978$, $\lambda_6 = 8/525 = 0.0152$ per year; logarithms $0.726$, $-0.119$, $-1.010$,
$-1.817$.

\textbf{22.} The line gives $10^{0.73} = 5.37$, $10^{-0.12} = 0.759$, $10^{-0.97} =
0.107$ and $10^{-1.82} = 0.0151$: within \qty{1}{\percent}, \qty{0}{\percent},
\qty{10}{\percent} and \qty{1}{\percent} of the counts.

\textbf{23.} The record of large eruptions stays complete further back in time,
because they are noticed from farther away and their deposits survive longer; a
window starting in 1500 would undercount VEI~3 eruptions badly, while a window
starting in 1950 would hold almost no VEI~6 eruption.

\textbf{24.} $\lambda_5 = \num{0.107}$ per year, a return period of $1/0.107 =
\num{9.3}\text{ years}$; $1 - e^{-1.07} = 1 - 0.343 = \qty{66}{\percent}$ in ten years.

\textbf{25.} $\lambda_7 = 10^{3.28 - 5.95} = 10^{-2.67} = 0.00214$ per year, one every
470 years; the record gives $2/1025 = 0.00195$ per year: agreement within
\qty{10}{\percent}, for a point not used in the fit.

\textbf{26.} $\lambda_6 = 10^{3.28 - 5.10} = 10^{-1.82} = 0.0151$ per year; $\lambda t = 1.51$
and $1 - e^{-1.51} = 1 - 0.220 = 0.780$: the probability of at least one eruption
of VEI~6 or more somewhere on Earth in the next hundred years is
\textbf{\qty[text-series-to-math=true]{78}{\percent}}. It rests on a constant rate,
on eruptions independent of one another (a Poisson process), and on a
catalogue complete within its windows.
\end{solution}
